\section{Experiments}
We evaluate \textsc{Dream-RSI} across three scientific discovery domains: algorithm engineering, kernel optimization and math optimization. Our primary controlled baseline is Recursive Fixed Exploration, which uses the same underlying discovery setting and initialization but keeps the exploration policy fixed across recursive discovery rounds. We additionally compare against task-specific domain baselines. 

Across all tasks, \textsc{Dream-RSI} and Recursive Fixed Exploration
use the same discovery agent, evaluator, initialization, and resource constraints.
Both methods start from the same manually designed exploration policy.
This exploration policy follows a simple \emph{parallel refining} strategy: it launches
multiple independent exploration workspaces in parallel, with each workspace
maintaining its own local discovery trajectory and repeatedly refining its
current candidate based on the history accumulated within that workspace.
The two methods therefore follow the same exploration policy in the first
discovery round. In subsequent rounds, while Recursive Fixed Exploration keeps its exploration policy static, \textsc{Dream-RSI} progressively refines the policy by dreaming over a replay simulator conditioned on accumulated global discovery history, subsequently deploying the updated policy in each new round. The discovery cost is quantified by the total cumulative number of discovery-agent calls.


Specifically, we evaluate Gemini-3.1 Pro and Gemini-3.7-Flash across multiple recursive discovery rounds via the Gemini CLI~\footnote{\url{https://geminicli.com/}}. Under Recursive Fixed Exploration, each round for Gemini-3.1 Pro executes 10 parallel workspaces with up to 11 refinement steps ($10 \times 11 = 110$ discovery-agent calls), whereas Gemini-3.7-Flash operates 32 parallel workspaces with up to 20 refinement steps ($32 \times 20 = 640$ calls). \textsc{Dream-RSI} maintains identical per-round budgets, aligning with the baseline in Round 1 while progressively updating its policy in subsequent rounds. Further details on recursive rounds, task setups, resource budgets, and evaluation protocols follow below.



\usepgfplotslibrary{groupplots}
\pgfplotsset{compat=1.18}

% ============================================================
% Scaling data
% Metric: arithmetic mean runtime over six real downstream tasks
% Lower is better.
% ============================================================

% ------------------------------------------------------------
% Gemini-3.1-Pro: Fixed Exploration
% ------------------------------------------------------------
\pgfplotstableread{
compute avg iter
110 5267.0500 1
220 5365.4000 2
330 5295.2333 3
440 4691.0833 4
550 3587.0667 5
}\probaseline

% ------------------------------------------------------------
% Gemini-3.1-Pro: Dream-RSI
% Iter 1 shares the same initialization
% ------------------------------------------------------------
\pgfplotstableread{
compute avg iter
110 5267.0500 1
119 5349.1500 2
147 5368.9000 3
234 5705.4333 4
317 2931.0000 5
}\proours

% ------------------------------------------------------------
% Gemini-3.7-Flash: Fixed Exploration
% Iter 1--2 have no downstream evaluation
% ------------------------------------------------------------
\pgfplotstableread{
compute avg iter
640  2735.6 1
1280 2988.3 2
1920 3011.1333 3
2560 2909.1833 4
3200 2516.6833 5
}\flashbaseline

% ------------------------------------------------------------
% Gemini-3.7-Flash: Dream-RSI
% Iter 1 has no downstream evaluation
% ------------------------------------------------------------
\pgfplotstableread{
compute avg iter
640  2735.6 1
976  2394.9167 2
1230 2378.0500 3
1529 2327.7000 4
1879 2350.5833 5
}\flashours


% ============================================================
% Main figure
% ============================================================

\begin{figure*}[t]
\centering

% ============================================================
% (a) Final performance table
% ============================================================

\resizebox{\textwidth}{!}{%
\begin{tabular}{l r r rr rrrr r}
\toprule

\textbf{Method}
& \textbf{Model}
& \textbf{Compute}
& \multicolumn{2}{c}{\textit{Non-biological}}
& \multicolumn{4}{c}{\textit{Biological}}
& \textbf{Avg.}
\\

\cmidrule(lr){4-5}
\cmidrule(lr){6-9}

&
&
&
\textbf{Gisette}
&
\textbf{RCV1}
&
\textbf{DNA}
&
\textbf{Leukemia}
&
\textbf{Colon}
&
\textbf{Duke Breast}
&
\\

\midrule

\multicolumn{10}{l}{\textit{Previous solvers}} \\[1pt]

sklearn
& --
& --
& 11275.2
& 252881.7
& 93.8
& 227.2
& 229.8
& 374.0
& 44180.3
\\

glmnet
& --
& --
& 9063.6
& 73072.8
& 351.9
& 45.0
& 24.2
& 47.7
& 13767.5
\\

SimpleTES
& gpt-oss-120b
& 51{,}200
& 3141.9
& 19625.6
& 15.9
& 15.5
& 11.6
& 18.1
& 3804.8
\\

SimpleTES $\dagger$
& gpt-oss-120b
& 51{,}200
& 8651.0
& 41143.1
& 37.6
& 28.2
& 19.5
& 31.1
& 8318.4
\\

\midrule

\multicolumn{10}{l}{\textit{Our System}} \\[1pt]

Recursive Fixed Exploration
& Gemini-3.1-Pro
& 550
& 1861.8
& 19550.1
& 41.5
& 26.1
& 14.5
& 28.4
& 3587.1
\\

&
Gemini-3.7-Flash
& 3200
& 1133.1
& 13873.0
& 29.8
& 24.1
& 15.7
& 24.4
& 2516.7
\\

\textbf{\textsc{Dream-RSI}}
& Gemini-3.1-Pro
& 317
& 2841.0
& 14616.0
& 49.9
& 30.2
& 16.4
& 32.5
& 2931.0
\\

&
Gemini-3.7-Flash
& 1879
& 1091.9
& 12923.4
& 31.4
& 21.0
& 12.2
& 23.6
& \textbf{2350.6}
\\

\bottomrule
\end{tabular}%
}

\vspace{0.08cm}

{\small\textbf{(a) Final performance.}}

\vspace{0.27cm}


% ============================================================
% (b) Scaling curves
% ============================================================

\begin{tikzpicture}

\begin{groupplot}[
group style={
  group size=2 by 1,
  horizontal sep=0.1\textwidth
},
width=0.5\textwidth,
height=4.85cm,
grid=major,
major grid style={
  gray!16,
  line width=0.3pt
},
axis line style={
  black!55,
  line width=0.65pt
},
tick style={
  black!55,
  line width=0.5pt
},
tick align=outside,
tick label style={
  font=\scriptsize,
  text=black!80
},
xlabel style={
  font=\small,
  yshift=1pt
},
ylabel style={
  font=\small
},
title style={
  font=\small\bfseries,
  yshift=2pt
},
scaled ticks=false,
clip=false
]


% ============================================================
% Left: Gemini-3.1-Pro
% ============================================================

\nextgroupplot[
title={Gemini-3.1-Pro},
xmin=50,
xmax=600,
ymin=2800,
ymax=6000,
xtick={100,300,500},
ytick={3000,4000,5000,6000},
xlabel={Discovery Compute},
ylabel={Avg. Runtime (ms)},
xticklabel style={
  /pgf/number format/fixed,
  /pgf/number format/precision=0
},
yticklabel style={
  /pgf/number format/fixed,
  /pgf/number format/precision=0,
  /pgf/number format/1000 sep={,}
},
legend to name=scalinglegend,
legend columns=2,
legend style={
  draw=none,
  font=\footnotesize,
  /tikz/every even column/.append style={
    column sep=1.4em
  }
}
]


% ------------------------------------------------------------
% Gemini-3.1-Pro: Fixed Exploration
% ------------------------------------------------------------

\addplot[
color=blue!75!black,
line width=1.7pt,
mark=square*,
mark size=2.5pt,
point meta=explicit symbolic,
nodes near coords,
every node near coord/.append style={
  font=\tiny,
  text=blue!75!black,
  anchor=north,
  yshift=-1.5pt,
  fill=white,
  fill opacity=0.90,
  text opacity=1,
  inner sep=0.7pt
}
]
table[
x=compute,
y=avg,
meta=iter
] {\probaseline};

\addlegendentry{Recursive Fixed Exploration}


% ------------------------------------------------------------
% Gemini-3.1-Pro: Dream-RSI
% ------------------------------------------------------------

\addplot[
color=red!80!black,
line width=1.9pt,
mark=*,
mark size=2.8pt,
point meta=explicit symbolic,
nodes near coords,
every node near coord/.append style={
  font=\tiny,
  text=red!80!black,
  anchor=south,
  yshift=1.5pt,
  fill=white,
  fill opacity=0.90,
  text opacity=1,
  inner sep=0.7pt
}
]
table[
x=compute,
y=avg,
meta=iter
] {\proours};

\addlegendentry{\textsc{Dream-RSI}}


% ============================================================
% Right: Gemini-3.7-Flash
% ============================================================

\nextgroupplot[
title={Gemini-3.7-Flash},
xmin=500,
xmax=3400,
ymin=2250,
ymax=3100,
xtick={1000,2000,3000},
ytick={2250,2500,2750,3000},
xlabel={Discovery Compute},
ylabel={},
xticklabel style={
  /pgf/number format/fixed,
  /pgf/number format/precision=0,
  /pgf/number format/1000 sep={,}
},
yticklabel style={
  /pgf/number format/fixed,
  /pgf/number format/precision=0,
  /pgf/number format/1000 sep={,}
}
]


% ------------------------------------------------------------
% Gemini-3.7-Flash: Fixed Exploration
% ------------------------------------------------------------

\addplot[
color=blue!75!black,
line width=1.7pt,
mark=square*,
mark size=2.5pt,
point meta=explicit symbolic,
nodes near coords,
every node near coord/.append style={
  font=\tiny,
  text=blue!75!black,
  anchor=south,
  yshift=1.5pt,
  fill=white,
  fill opacity=0.90,
  text opacity=1,
  inner sep=0.7pt
}
]
table[
x=compute,
y=avg,
meta=iter
] {\flashbaseline};


% ------------------------------------------------------------
% Gemini-3.7-Flash: Dream-RSI
% ------------------------------------------------------------

\addplot[
color=red!80!black,
line width=1.9pt,
mark=*,
mark size=2.8pt,
point meta=explicit symbolic,
nodes near coords,
every node near coord/.append style={
  font=\tiny,
  text=red!80!black,
  anchor=south,
  yshift=1.5pt,
  fill=white,
  fill opacity=0.90,
  text opacity=1,
  inner sep=0.7pt
}
]
table[
x=compute,
y=avg,
meta=iter
] {\flashours};


\end{groupplot}

\end{tikzpicture}


% ============================================================
% Shared legend
% ============================================================

\vspace{0.00cm}

{\centering
\pgfplotslegendfromname{scalinglegend}
\par
}

\vspace{0.015cm}

{\small\textbf{(b) Recursive Discovery Dynamics.}}

\vspace{-0.05cm}


% ============================================================
% Caption
% ============================================================

\caption{
\textbf{Lasso regularization-path discovery results.}
\textbf{(a)} Final wall-clock runtime on six held-out downstream tasks;
lower is better.
\textbf{Compute} denotes the cumulative number of discovery-agent calls.
\textbf{(b)} \textbf{Recursive discovery dynamics.} Average downstream runtime across six held-out tasks versus cumulative discovery compute for Gemini-3.1-Pro and Gemini-3.7-Flash. Numbers next to markers denote recursive rounds (iterations). Lower is better.
}

\label{fig:lasso-main}

\end{figure*}


\subsection{Algorithm Engineering}
\label{subsec:lasso}

In this task, we consider \textbf{Lasso Regularization Path} as our algorithm-engineering task, a fundamental computational primitive in high-dimensional statistics that is widely used in model selection and cross-validation across domains such as genomics and finance. We follow the benchmark setting of SimpleTES~\citep{ye2026evaluation}, where the goal is to discover efficient implementations of the complete Lasso regularization path while preserving numerical correctness. During discovery, we use the same 17 synthetic instances as SimpleTES, which cover diverse problem regimes in terms of dimensionality, sparsity, feature correlation, and active-set structure. To evaluate whether the discovered algorithms generalize beyond the search distribution, we additionally evaluate them on six held-out downstream datasets spanning both biological and non-biological domains. 

\paragraph{Baselines and Setup.}
We compare against standard Lasso solvers
\textbf{sklearn}~\citep{pedregosa2011scikit} and
\textbf{glmnet}~\citep{friedman2010regularization},
as well as \textbf{SimpleTES}~\citep{ye2026evaluation}, which uses
GPT-OSS-120B with a reported budget of 51,200 generations.
We additionally include Recursive Fixed Exploration as our controlled baseline. Specifically, we run both Recursive Fixed Exploration and \textsc{Dream-RSI} for 5 rounds.

\paragraph{Main Results.}
Figure~\ref{fig:lasso-main}(a) summarizes the Lasso discovery results.
Across both discovery-agent backbones, \textsc{Dream-RSI} achieves a better
downstream quality--compute trade-off than Recursive Fixed Exploration.
With Gemini-3.1 Pro, it reduces the average runtime across the six held-out
datasets from 3587.1\,ms to 2931.0\,ms while using only 317 discovery-agent
calls, compared with 550 calls for fixed exploration.
With Gemini-3.7-Flash, \textsc{Dream-RSI} further reduces the average runtime
from 2516.7\,ms to 2350.6\,ms using 1879 calls instead of 3200. Despite using substantially less discovery compute, the resulting solvers also
outperform the standard sklearn and glmnet implementations on all six held-out
datasets. Compared with SimpleTES, which uses 51,200 generations,
\textsc{Dream-RSI} achieves lower average downstream runtime with roughly
two orders of magnitude fewer discovery-agent calls. Notably, the program discovered by Gemini-3.1-Pro appears particularly well suited to large-scale matrices such as RCV1. In contrast, Gemini-3.7-Flash discovers a more general-purpose program that performs consistently across different problem scales.

\paragraph{Recursive Discovery Dynamics.}
Figure~\ref{fig:lasso-main}(b) illustrates the trajectory of downstream performance across recursive discovery rounds relative to cumulative discovery compute. By design, both methods share identical search behavior in the initial round. In subsequent rounds, Recursive Fixed Exploration maintains a static exploration policy, whereas \textsc{Dream-RSI} progressively refines and redeploys its policy via dreaming over accumulated discovery history. Consequently, the two trajectories diverge markedly: \textsc{Dream-RSI} consistently achieves superior downstream performance while requiring substantially lower cumulative compute across both Gemini-3.1-Pro and Gemini-3.7-Flash.


\paragraph{Discovered Solver Analysis.}
We further analyze the discovered solver, with its implementation provided in
the Appendix \ref{app:discovered-lasso}. Unlike SimpleTES, which switches between LARS and coordinate
descent according to problem dimensions, the discovered solver introduces
adaptivity within the active-set optimization itself. It combines strong-rule
screening with Cauchy--Schwarz-based KKT pruning, selectively recomputing exact
gradients only when the bound cannot certify a feature and falling back to a
full refresh when pruning becomes ineffective. This adaptive verification
scheme is further integrated with efficient active-set bookkeeping, lazy
Gram-matrix construction, and hardware-aware implementation.


\begin{table}[t!]
\centering
\caption{
Performance comparison on mathematical discovery tasks.
Higher is better for \textbf{Sum Diff} and \textbf{Circle Packing}, while lower is better for \textbf{Auto Correlation}.
Best results are shown in \textbf{bold}.
}
\label{tab:downstream_results}
\small
\setlength{\tabcolsep}{1pt}
\renewcommand{\arraystretch}{1.15}
\vspace{0.1in}
\begin{tabular}{llccc}
\toprule
\textbf{Method}
& \textbf{LLM}
& \textbf{Sum Diff ($\uparrow$)}
& \textbf{Auto Correlation ($\downarrow$)}
& \textbf{Circle Packing ($\uparrow$)} \\
\midrule

AlphaEvolve
& Gemini-2.0 Pro + Flash 
& -- &  1.455700 &  2.635862 \\

AlphaEvolveV2
& Gemini-2.0 Pro + Flash 
& 1.121936 & -- & \textbf{2.635983} \\

OpenEvolve 
& - 
& -- &  1.460000 &  - \\

CodeEvolve 
& -
& -- & -- & 2.635980 \\

ShinkaEvolve 
& Mixed 
& -- &  1.457800 &  2.635982 \\

TTS-Discovery
& Qwen3-8B 
& -- & -- &  \textbf{2.635983} \\

ThetaEvolve
& Distilled-Qwen3-8B 
& -- & 1.493000 & \textbf{2.635983} \\

EvoX
& Gemini-3.0-Pro 
& -- & 1.458900 & 2.635900 \\

SimpleTES
& GPT-OSS-120B 
& 1.143975 & \bf 1.453675 & \textbf{2.635983} \\


\midrule
\textit{Our System} \\
Recursive Fixed Exploration
& Gemini-3.1-Pro 
& 1.144047 & 1.456001 & \textbf{2.635983} \\

\textsc{Dream-RSI}
& Gemini-3.1-Pro  
& \textbf{1.145427}
& 1.456375
& \textbf{2.635983} \\

\bottomrule
\end{tabular}
\end{table}



\subsection{Mathematics Optimization}

We further evaluate \textsc{Dream-RSI} on three mathematical discovery tasks spanning discrete combinatorial optimization, geometric optimization, and functional optimization: the \textbf{Sum--Difference Problem}, \textbf{Circle Packing}, and \textbf{Autocorrelation Inequalities}. The goal of these problems is to discover high-quality solutions that optimize task-specific mathematical objectives under their respective constraints. Formal definitions of the three tasks are provided in Appendix.

We use Gemini-3.1 Pro via the Gemini CLI as the discovery agent for both Recursive Fixed Exploration and \textsc{Dream-RSI} for 10 rounds. For each task, the agent iteratively proposes and evaluates candidate constructions or optimization procedures according to the task-specific objective. We compare against a broad set of existing automated discovery systems, including AlphaEvolve~\citep{novikov2025alphaevolve}, AlphaEvolveV2~\citep{georgiev2025mathematical}, OpenEvolve~\citep{openevolve}, CodeEvolve~\citep{assumpccao2025codeevolve}, ShinkaEvolve~\citep{lange2026shinkaevolve}, TTS-Discovery~\citep{yuksekgonul2026learning}, ThetaEvolve~\citep{wang2025thetaevolve}, EvoX~\citep{liu2026evox}, and SimpleTES~\citep{ye2026evaluation}.

\paragraph{Results.}
Table~\ref{tab:downstream_results} summarizes the results across the three mathematical discovery tasks.
\textsc{Dream-RSI} achieves a Sum--Difference score of $1.145427$, outperforming SimpleTES and Recursive Fixed Exploration.
On Circle Packing, it reaches $2.635983$, matching the strongest reported result among the compared methods.
For Autocorrelation, \textsc{Dream-RSI} obtains $1.456375$, remaining competitive with existing discovery systems. Notably, SimpleTES achieves state-of-the-art performance on Autocorrelation Inequalities, but requires 51,200 generations, significantly more than the fewer than 1,000 generations used by our approach.
Overall, these results show that our \textsc{Dream-RSI} generalize well on mathematics optimization.



\begin{figure}[t]
    \centering
    \definecolor{oursblue}{HTML}{0072B2}
\definecolor{baselinegray}{HTML}{6B7280}
\usepgfplotslibrary{groupplots}
\begin{tikzpicture}

\begin{groupplot}[
  group style={group size=2 by 2,horizontal sep=1.00cm,vertical sep=1.45cm},
  width=7.3cm,
  height=4.4cm,
  xmin=0,
  axis lines=left,
  axis line style={black!55,line width=0.65pt},
  xmajorgrids=false,
  ymajorgrids=true,
  grid style={black!8,line width=0.30pt},
  scaled ticks=false,
  tick align=outside,
  tick style={black!55,line width=0.50pt},
  tick label style={font=\scriptsize,text=black!80},
  title style={font=\small\bfseries,yshift=3pt},
  xlabel style={font=\small,yshift=2pt},
  ylabel style={font=\small},
  clip=true,
  clip marker paths=true,
  restrict x to domain=1:2000,
  unbounded coords=discard,
  filter discard warning=false
]

% ============================================================
% VGG16
% ============================================================

\nextgroupplot[
  title={VGG16},
  xmax=1050,
  ymin=0.35,
  ymax=0.57,
  xtick={0,250,500,750,1000},
  xticklabels=\empty,
  ytick={0.35,0.40,0.45,0.50,0.55},
  ylabel={Performance ($1/\mathrm{ms}$)},
  legend to name=sharedlegend,
  legend columns=2,
  legend style={draw=none,font=\small,column sep=1.4em}
]

\addplot[
  const plot,
  color=oursblue,
  line width=2.55pt,
  mark=*,
  mark size=1.35pt,
  mark options={fill=oursblue,draw=oursblue}
]
table[x=budget,y=score]
{Main_Text/raw_data/data/vgg16_ours.dat};

\addlegendentry{\textsc{Dream-RSI}}

\addplot[
  const plot,
  color=baselinegray,
  line width=1.75pt,
  densely dashed,
  mark=square*,
  mark size=1.00pt,
  mark options={fill=white,draw=baselinegray,line width=0.65pt}
]
table[x=budget,y=score]
{Main_Text/raw_data/data/vgg16_baseline.dat};

\addlegendentry{Recursive Fixed Exploration}

\node[
  anchor=north west,
  font=\scriptsize\bfseries,
  text=oursblue!80!black,
  fill=white,
  draw=oursblue!20,
  rounded corners=1pt,
  line width=0.3pt,
  inner xsep=3pt,
  inner ysep=1.7pt
]
at (axis description cs:0.38,0.67)
{\makecell{$2.43\times$ fewer costs \\ (comparable performance)}};


% ============================================================
% LayerNorm
% ============================================================

\nextgroupplot[
  title={LayerNorm},
  xmax=1050,
  ymin=1.02,
  ymax=1.15,
  xtick={0,250,500,750,1000},
  xticklabels=\empty,
  ytick={1.00,1.05,1.10,1.15}
]

\addplot[
  const plot,
  color=oursblue,
  line width=2.55pt,
  mark=*,
  mark size=1.35pt,
  mark options={fill=oursblue,draw=oursblue}
]
table[x=budget,y=score]
{Main_Text/raw_data/data/layernorm_ours.dat};

\addplot[
  const plot,
  color=baselinegray,
  line width=1.75pt,
  densely dashed,
  mark=square*,
  mark size=1.00pt,
  mark options={fill=white,draw=baselinegray,line width=0.65pt}
]
table[x=budget,y=score]
{Main_Text/raw_data/data/layernorm_baseline.dat};

\node[
  anchor=north west,
  font=\scriptsize\bfseries,
  text=oursblue!80!black,
  fill=white,
  draw=oursblue!20,
  rounded corners=1pt,
  line width=0.3pt,
  inner xsep=3pt,
  inner ysep=1.7pt
]
at (axis description cs:0.335,0.55)
{\makecell{$1.79\times$ fewer costs \\ (comparable performance)}};


% ============================================================
% ConvDiv
% ============================================================

\nextgroupplot[
  title={ConvDiv},
  xmax=1050,
  ymin=0.40,
  ymax=2.00,
  xtick={0,250,500,750,1000},
  ytick={0.4,0.8,1.2,1.6,2.0},
  xlabel={Number of Generations},
  ylabel={Performance ($1/\mathrm{ms}$)}
]

\addplot[
  const plot,
  color=oursblue,
  line width=2.55pt,
  mark=*,
  mark size=1.35pt,
  mark options={fill=oursblue,draw=oursblue}
]
table[x=budget,y=score]
{Main_Text/raw_data/data/convdiv_ours.dat};

\addplot[
  const plot,
  color=baselinegray,
  line width=1.75pt,
  densely dashed,
  mark=square*,
  mark size=1.00pt,
  mark options={fill=white,draw=baselinegray,line width=0.65pt}
]
table[x=budget,y=score]
{Main_Text/raw_data/data/convdiv_baseline.dat};

\node[
  anchor=north west,
  font=\scriptsize\bfseries,
  text=oursblue!80!black,
  fill=white,
  draw=oursblue!20,
  rounded corners=1pt,
  line width=0.3pt,
  inner xsep=3pt,
  inner ysep=1.7pt
]
at (axis description cs:0.035,0.95)
{\makecell{$2.09\times$ higher score \\ (Similar Budget)}};


% ============================================================
% ConvMax
% ============================================================

\nextgroupplot[
  title={ConvMax},
  xmax=1050,
  ymin=0.22,
  ymax=0.45,
  xtick={0,250,500,750,1000},
  ytick={0.25,0.30,0.35,0.40,0.45},
  xlabel={Number of Generations}
]

\addplot[
  const plot,
  color=oursblue,
  line width=2.55pt,
  mark=*,
  mark size=1.35pt,
  mark options={fill=oursblue,draw=oursblue}
]
table[x=budget,y=score]
{Main_Text/raw_data/data/convmax_ours.dat};

\addplot[
  const plot,
  color=baselinegray,
  line width=1.75pt,
  densely dashed,
  mark=square*,
  mark size=1.00pt,
  mark options={fill=white,draw=baselinegray,line width=0.65pt}
]
table[x=budget,y=score]
{Main_Text/raw_data/data/convmax_baseline.dat};

\node[
  anchor=north west,
  font=\scriptsize\bfseries,
  text=oursblue!80!black,
  fill=white,
  draw=oursblue!20,
  rounded corners=1pt,
  line width=0.3pt,
  inner xsep=3pt,
  inner ysep=1.7pt
]
at (axis description cs:0.035,0.95)
{\makecell{$1.44\times$ higher score \\ (Similar Budget)}};

\end{groupplot}

% ============================================================
% Shared legend on top
% ============================================================

\node[anchor=south]
at ([yshift=0.05cm]current bounding box.north)
{\pgfplotslegendfromname{sharedlegend}};

\end{tikzpicture}
    \caption{
\textbf{GPU kernel engineering results.}
Discovery performance of \textsc{Dream-RSI} and Recursive Fixed Exploration as a function of the number of generations.
On VGG16 and LayerNorm, \textsc{Dream-RSI} reaches comparable performance with $2.43\times$ and $1.79\times$ fewer generations, respectively.
On ConvDiv and ConvMax, it achieves $2.09\times$ and $1.44\times$ higher performance under comparable discovery budgets.
Higher is better for all tasks.
}
    \label{fig:kernelbench-scaling}
\end{figure}
\subsection{Kernel Engineering}

We further evaluate \textsc{Dream-RSI} on \textbf{GPU kernel engineering}, where the goal is to automatically discover high-performance implementations of kernels while preserving numerical correctness. Unlike mathematical discovery, kernel engineering requires reasoning jointly about algorithmic structure, memory access, parallelization, and hardware-specific optimizations, providing a substantially different testbed for evaluating whether our \textsc{Dream-RSI} generalizes across discovery domains.

We consider four representative kernel-engineering tasks from KernelBench~\citep{ouyang2025kernelbench}: \textbf{VGG16}, \textbf{LayerNorm}, \textbf{ConvDiv}, and \textbf{ConvMax}. Candidate implementations are evaluated by their execution performance, measured as inverse runtime ($1/\mathrm{ms}$), subject to correctness checks against the reference implementation. We use Gemini-3.1 Pro as the coding agent and compare \textsc{Dream-RSI} with Recursive Fixed Exploration under the same evaluation protocol and initialization.

\paragraph{Results.}
Figure~\ref{fig:kernelbench-scaling} shows the discovery trajectories as the number of generations increases.
On VGG16 and LayerNorm, \textsc{Dream-RSI} reaches comparable final performance using
$2.43\times$ and $1.79\times$ fewer generations, respectively.
On ConvDiv and ConvMax, under comparable discovery budgets,
\textsc{Dream-RSI} achieves $2.09\times$ and $1.44\times$ higher performance, respectively.
These results show that adapting the exploration policy across recursive rounds can improve the efficiency and effectiveness of long-horizon discovery.





